\documentclass[12pt]{article}
\usepackage[utf8]{inputenc}
\usepackage[T1]{fontenc}
\usepackage[margin=1in]{geometry}
\usepackage{setspace}
\usepackage{xcolor}
\usepackage[colorlinks=true, urlcolor=blue!60!black, linkcolor=blue!60!black]{hyperref}
\setlength{\parindent}{2em}
\setlength{\parskip}{1.25em}
\renewcommand{\baselinestretch}{1.0}
\begin{document}
\begin{center}
{\LARGE \textit{Fletcher v. Peck}, 10 U.S. 87 (1810)\par}
\vspace{0.8em}
{\large Case Brief\par}
\end{center}
\vspace{1em}

\section{Memory Jogger}

In the first Supreme Court decision invalidating a state statute as unconstitutional, Chief Justice Marshall held that a state's legislative land grant is an executed ``contract'' under Article~I, \S~10's Contract Clause, barring Georgia from rescinding its 1795 Yazoo land grant against innocent third-party purchasers even though the original statute was procured by massive legislative bribery (\href{https://www.courtlistener.com/opinion/84935/fletcher-v-peck/}{CourtListener}; \href{https://www.fjc.gov/history/cases/cases-that-shaped-the-federal-courts/fletcher-v-peck}{Federal Judicial Center}).

\section{Detailed Case Facts}

On January 7, 1795, the Georgia General Assembly enacted a statute authorizing the sale of approximately 35 million acres of western territory along the Yazoo River (comprising much of present-day Alabama and Mississippi) to four private land companies, including the Georgia Company led by James Gunn, Matthew McAllister, and George Walker, for roughly 1.5 cents per acre (\href{https://supreme.justia.com/cases/federal/us/10/87/}{Justia}; \href{https://ballotpedia.org/Fletcher_v._Peck}{Ballotpedia}). The sale was procured through pervasive corruption: land-company representatives promised numerous Georgia legislators shares in the land if they voted for the bill, and the bribed legislators passed the act (\href{https://www.oyez.org/cases/1789-1850/10us87}{Oyez}; \href{https://supreme.justia.com/cases/federal/us/10/87/}{Justia}).

Public outrage over the Yazoo land fraud swept most of the implicated legislators out of office in the ensuing election. In 1796, the newly elected Georgia legislature enacted a rescinding statute declaring the 1795 grant null and void, asserting state title to the lands, and ordering all public records of the 1795 act expunged (\href{https://www.oyez.org/cases/1789-1850/10us87}{Oyez}). Before the 1796 repeal, however, the original grantees had already resold large tracts to northern speculators and innocent third-party purchasers for value without notice of the bribery, including members of the New England Mississippi Land Company (\href{https://www.fjc.gov/history/cases/cases-that-shaped-the-federal-courts/fletcher-v-peck}{Federal Judicial Center}; \href{https://www.americanhistorycentral.com/entries/fletcher-v-peck-1810/}{American History Central}).

John Peck, a Massachusetts land speculator who acquired title to a portion of the Yazoo lands in 1800, conveyed 15,000 acres to Robert Fletcher of New Hampshire on May 14, 1803, for \$3,000 (\href{https://www.oyez.org/cases/1789-1850/10us87}{Oyez}; \href{https://www.americanhistorycentral.com/entries/fletcher-v-peck-1810/}{American History Central}). Peck's deed contained four express covenants: (1) that the Georgia legislature had authority to sell and dispose of the lands; (2) that the 1795 act was valid and lawfully passed; (3) that Georgia's title had not been constitutionally or legally impaired by any subsequent legislative act (namely, the 1796 rescinding statute); and (4) that the State of Georgia was legally seised in fee simple of the lands at the time of the 1795 grant, despite unextinguished Native American occupancy (\href{https://www.oyez.org/cases/1789-1850/10us87}{Oyez}; \href{https://press-pubs.uchicago.edu/founders/documents/a1_10_1s12.html}{Founders' Constitution}). Because the Eleventh Amendment barred out-of-state landholders from suing Georgia directly in federal court, Fletcher and Peck arranged what historians and Justice William Johnson recognized as a ``friendly'' or ``feigned'' test case to obtain a federal judicial ruling on the validity of the Yazoo titles (\href{https://www.fjc.gov/history/cases/cases-that-shaped-the-federal-courts/fletcher-v-peck}{Federal Judicial Center}; \href{https://www.americanhistorycentral.com/entries/fletcher-v-peck-1810/}{American History Central}).

\section{Procedural History}

Four days after the May 14, 1803 conveyance, Fletcher sued Peck for breach of covenant in the U.S. Circuit Court for the District of Massachusetts, invoking diversity jurisdiction (\href{https://www.americanhistorycentral.com/entries/fletcher-v-peck-1810/}{American History Central}; \href{https://www.oyez.org/cases/1789-1850/10us87}{Oyez}). Fletcher's declaration pleaded four counts corresponding to the four covenants in Peck's deed: that Georgia lacked authority to convey the land under its 1789 state constitution (Count 1); that legislative bribery rendered the 1795 grant void \textit{ab initio} (Count 2); that Georgia's 1796 rescinding act validly extinguished the title (Count 3); and that Georgia was not seised in fee simple because the land lay outside its boundaries or remained subject to Indian title (Count 4) (\href{https://supreme.justia.com/cases/federal/us/10/87/}{Justia}; \href{https://www.law.cornell.edu/supremecourt/text/10/87}{LII}).

The parties postponed litigation in the circuit court for several years while Congress considered a federal compensation bill for Yazoo claimants (\href{https://www.fjc.gov/history/cases/cases-that-shaped-the-federal-courts/fletcher-v-peck}{Federal Judicial Center}). When legislative relief stalled, the parties resumed the suit. Peck pleaded that he and intervening purchasers were bona fide purchasers for valuable consideration without notice of the alleged fraud, and that Georgia was seised in fee simple; Fletcher filed general demurrers to Peck's pleas on the first three counts, and a jury returned a special verdict on the fourth count (\href{https://supreme.justia.com/cases/federal/us/10/87/}{Justia}). The Circuit Court rendered judgment for Peck on all four counts, upholding the 1795 grant and finding the 1796 rescinding act unconstitutional (\href{https://www.oyez.org/cases/1789-1850/10us87}{Oyez}; \href{https://www.fjc.gov/history/cases/cases-that-shaped-the-federal-courts/fletcher-v-peck}{Federal Judicial Center}).

Fletcher brought a writ of error to the U.S. Supreme Court, where the case was argued on March 1--4, 1809, and, following a technical pleading amendment, reargued on February 15, 1810 (\href{https://supreme.justia.com/cases/federal/us/10/87/}{Justia}; \href{https://ballotpedia.org/Fletcher_v._Peck}{Ballotpedia}). On March 16, 1810, the Supreme Court affirmed the judgment of the Circuit Court (\href{https://supreme.justia.com/cases/federal/us/10/87/}{Justia}).

\section{Judicial Votes}

The Court decided the case by a vote of \textbf{4--1} (or unanimous 5--0 as to the judgment affirming Peck's title on the rescission issue, with Justice Johnson disagreeing on the Contract Clause rationale and dissenting on the Indian-title fee-simple issue), with five of the seven sitting Justices participating (\href{https://www.supremecourt.gov/pdfs/USReports/USREPORTS-10_PDFA.pdf}{10 U.S. Reports (1810)}; \href{https://ballotpedia.org/Fletcher_v._Peck}{Ballotpedia}):

\begin{itemize}
\item \textbf{Majority (4):} Chief Justice John Marshall (author of the opinion of the Court), joined by Associate Justices Bushrod Washington, Henry Brockholst Livingston, and Thomas Todd (\href{https://ballotpedia.org/Fletcher_v._Peck}{Ballotpedia}).
\item \textbf{Concurring in Part and Dissenting in Part (1):} Associate Justice William Johnson, Jr., who agreed that Georgia could not revoke its land grant (resting on general principles of natural justice rather than the Contract Clause) but dissented from the majority's Contract Clause interpretation and from its conclusion that Georgia held a fee-simple estate prior to extinguishment of Indian title (\href{https://www.oyez.org/cases/1789-1850/10us87}{Oyez}; \href{https://press-pubs.uchicago.edu/founders/documents/a1_10_1s12.html}{Founders' Constitution}).
\item \textbf{Did Not Participate (2):} Associate Justices William Cushing and Samuel Chase, both of whom were prevented from attending the February 1810 Term by ill health (\href{https://www.supremecourt.gov/pdfs/USReports/USREPORTS-10_PDFA.pdf}{10 U.S. Reports (1810)}).
\end{itemize}

\section{Holding}

\paragraph{Judgment.}
Judgment of the U.S. Circuit Court for the District of Massachusetts \textbf{affirmed} in favor of defendant John Peck (\href{https://www.courtlistener.com/opinion/84935/fletcher-v-peck/}{CourtListener}; \href{https://ballotpedia.org/Fletcher_v._Peck}{Ballotpedia}).

\paragraph{Rule Announced.}
Under Article~I, \S~10, Clause~1 of the U.S. Constitution (the Contract Clause), a completed legislative land grant by a state is an executed contract whose obligations continue to bind the grantor state; once legal title has vested and passed into the hands of bona fide third-party purchasers for valuable consideration without notice, a subsequent state legislature cannot constitutionally repeal or rescind the grant, nor may a court invalidate an otherwise properly enacted statute in a collateral suit between private parties by inquiring into alleged corrupt motives of individual legislators (\href{https://www.courtlistener.com/opinion/84935/fletcher-v-peck/}{CourtListener}; \href{https://constitutionallawreporter.com/2015/06/16/fletcher-v-peck-contract-clause/}{Constitutional Law Reporter}). Furthermore, a state holds a fee-simple estate in lands within its borders capable of being conveyed even while subject to unextinguished Native American right of occupancy (\href{https://supreme.justia.com/cases/federal/us/10/87/}{Justia}).

\paragraph{Holding vs. Dicta.}
The \textbf{holding} encompasses four points necessary to affirm judgment for Peck on all four counts: (1) the 1795 Georgia statute was not void under Georgia's 1789 constitution; (2) a federal court in a private breach-of-covenant suit will not void a duly enacted state statute by probing legislators' corrupt motives; (3) Georgia's 1796 rescinding act was unconstitutional as applied to innocent third-party purchasers because a state land grant is an executed contract protected by the Contract Clause; and (4) Georgia's fee-simple seisin was not incompatible with unextinguished Indian title (\href{https://supreme.justia.com/cases/federal/us/10/87/}{Justia}; \href{https://www.oyez.org/cases/1789-1850/10us87}{Oyez}). By contrast, Chief Justice Marshall's broad remarks suggesting that Georgia might also be restrained by unwritten ``general principles which are common to our free institutions'' and ``the nature of society and of government''---as well as his analogy comparing the rescinding act to a bill of attainder or a civil \textit{ex post facto} law (despite \textit{Calder v. Bull}, 3 U.S. 386 (1798), limiting the Ex Post Facto Clause to criminal laws)---are \textbf{dicta}, because the Contract Clause independently supplied the dispositive federal constitutional bar (\href{https://www.fjc.gov/history/cases/cases-that-shaped-the-federal-courts/fletcher-v-peck}{Federal Judicial Center}; \href{https://supreme.justia.com/cases/federal/us/349/302/}{Justia}). Likewise, Marshall's suggestion that a state might possibly seek equitable rescission directly against original corrupt grantees who still held the land was unnecessary to decide Peck's case (\href{https://www.law.cornell.edu/supremecourt/text/10/87}{LII}).

\section{Analysis of Opinions}

\subsection{Majority Opinion (Chief Justice John Marshall)}

Chief Justice Marshall structured the Court's opinion sequentially across the four counts in Fletcher's declaration, establishing federal judicial review over state statutes in the process (\href{https://www.fjc.gov/history/cases/cases-that-shaped-the-federal-courts/fletcher-v-peck}{Federal Judicial Center}). First, addressing judicial restraint, Marshall cautioned that declaring a law void for repugnancy to a constitution is ``at all times, a question of much delicacy'' requiring clear incompatibility, and concluded that nothing in Georgia's 1789 constitution barred the legislature from selling unappropriated public lands (\href{https://caselaw.findlaw.com/court/us-supreme-court/10/87.html}{FindLaw}).

Second, on the bribery allegation, Marshall declined to permit a collateral attack on a duly enacted statute through a private suit between individuals. He reasoned that courts cannot feasibly investigate the impure motives of individual legislators or determine how many tainted votes would suffice to nullify a statute, particularly when the State of Georgia was not even a party to the record (\href{https://www.fjc.gov/history/cases/cases-that-shaped-the-federal-courts/fletcher-v-peck}{Federal Judicial Center}). Even if a court of equity could unwind a fraudulent conveyance between original wrongdoers, settled principles of equity protect bona fide purchasers for valuable consideration without notice (\href{https://www.oyez.org/cases/1789-1850/10us87}{Oyez}).

Third, turning to the 1796 rescinding act, Marshall distinguished ordinary repealable legislation from completed transfers of vested rights. While one legislature may ordinarily repeal a prior legislature's general statutes, ``when a law is in its nature a contract, when absolute rights have vested under that contract, a repeal of the law cannot devest those rights'' (\href{https://press-pubs.uchicago.edu/founders/documents/a1_10_1s12.html}{Founders' Constitution}). Marshall then anchored this limit in Article~I, \S~10's Contract Clause. He defined a contract as a compact that is either executory (binding a party to do or not do a future act) or executed (performing the object of the contract immediately), reasoning that a grant is an executed contract that extinguishes the grantor's right and implies a continuing obligation not to reassert it (\href{https://ballotpedia.org/Fletcher_v._Peck}{Ballotpedia}). Because the constitutional text forbids any state from passing any law impairing the obligation of contracts without distinguishing between private contracts and public contracts, the Clause binds a state's own grants just as it protects private agreements (\href{https://constitutionallawreporter.com/2015/06/16/fletcher-v-peck-contract-clause/}{Constitutional Law Reporter}). Reinforcing this textual reading, Marshall observed that the Framers viewed state legislative excesses with apprehension and designed Article~I, \S~10 as a ``bill of rights for the people of each state,'' comparing Georgia's confiscation to a bill of attainder and an \textit{ex post facto} law (\href{https://caselaw.findlaw.com/court/us-supreme-court/10/87.html}{FindLaw}; \href{https://www.fjc.gov/history/cases/cases-that-shaped-the-federal-courts/fletcher-v-peck}{Federal Judicial Center}). He concluded that Georgia was restrained ``either by general principles which are common to our free institutions, or by the particular provisions of the constitution of the United States'' (\href{https://constitutionallawreporter.com/2015/06/16/fletcher-v-peck-contract-clause/}{Constitutional Law Reporter}).

Fourth, on the special verdict regarding seisin, Marshall held that the Yazoo lands lay within Georgia's borders and that the unextinguished Indian title---which he characterized as a temporary right of occupancy to be respected until lawfully extinguished---was not incompatible with Georgia's holding a fee-simple estate (\href{https://supreme.justia.com/cases/federal/us/10/87/}{Justia}).

\subsection{Separate Opinion: Concurring in Part and Dissenting in Part (Justice William Johnson)}

Justice Johnson agreed that Georgia lacked power to revoke its own land grant, but he sharply rejected Marshall's reliance on the Contract Clause (\href{https://www.oyez.org/cases/1789-1850/10us87}{Oyez}; \href{https://press-pubs.uchicago.edu/founders/documents/a1_10_1s12.html}{Founders' Constitution}). Instead, Johnson grounded the invalidity of the 1796 rescinding act strictly in natural law---``on a general principle, on the reason and nature of things: a principle which will impose laws even on the deity'' (\href{https://caselaw.findlaw.com/court/us-supreme-court/10/87.html}{FindLaw}). Once a state grants property, the right passes completely out of the state; attempting to take it back by legislative fiat is an act of naked power rather than law.

Johnson warned that stretching the Contract Clause to cover executed public grants went ``very far beyond the obvious and necessary import of the words'' and risked paralyzing ordinary state legislation, including statutes of limitations, recording acts, and the essential sovereign power of eminent domain (\href{https://press-pubs.uchicago.edu/founders/documents/a1_10_1s12.html}{Founders' Constitution}). He also dissented from the majority's holding on the fourth count, arguing that while Native American nations held an unextinguished right of occupancy, Georgia could not be seised in fee simple of the land, holding at most a future interest contingent on extinguishment of Indian title (\href{https://press-pubs.uchicago.edu/founders/documents/a1_10_1s12.html}{Founders' Constitution}). Finally, Johnson openly voiced misgivings about hearing the case at all, observing that it bore ``strong evidence, upon the face of it, of being a mere feigned case,'' and proceeded only out of respect for the eminent counsel arguing the cause (\href{https://www.uclalawreview.org/wp-content/uploads/2019/09/12_52UCLALRev2892004-2005.pdf}{UCLA Law Review}; \href{https://www.fjc.gov/history/cases/cases-that-shaped-the-federal-courts/fletcher-v-peck}{Federal Judicial Center}).

\section{Examples: Future Applications}

\begin{enumerate}
\item \textbf{Same-Side (Unconstitutional Impairment of Public Grant):} A state legislature passes a statute granting a private university a permanent corporate charter and tax-exempt endowment. Ten years later, a newly elected legislature amends the charter to convert the university into a state-run public institution without the trustees' consent. Under \textit{Fletcher v. Peck}, the repeal is unconstitutional because a state-granted charter, like a land grant, is an executed contract whose vested rights cannot be unilaterally revoked by a subsequent legislature (mirroring \textit{Trustees of Dartmouth College v. Woodward}, 17 U.S. 518 (1819)) (\href{https://www.fjc.gov/history/cases/cases-that-shaped-the-federal-courts/fletcher-v-peck}{Federal Judicial Center}; \href{https://press-pubs.uchicago.edu/founders/documents/a1_10_1s17.html}{Founders' Constitution}).

\item \textbf{Same-Side (Repeal of Curative Tax-Sale Statute):} A state sells tax-forfeited real estate to private investors under a statutory warranty that cures minor procedural irregularities in the tax foreclosure. After a political backlash against investor purchases, the legislature repeals the curative statute retroactively to strip title from bona fide buyers who already purchased and resold the parcels. Under \textit{Fletcher}, the retroactive repeal violates the Contract Clause because it destroys vested property rights conveyed by the state to purchasers for value (mirroring \textit{Wood v. Lovett}, 313 U.S. 362 (1941)) (\href{https://supreme.justia.com/cases/federal/us/313/362/}{Justia}).

\item \textbf{Opposite-Side (Prospective Repeal of Unvested Statutory Benefit):} A state statute offers a renewable annual subsidy to solar-panel manufacturers who build factories in the state, without entering into binding multi-year contracts or conveying title to any property. Three years later, facing a budget shortfall, the legislature repeals the subsidy program for all future tax years. Unlike \textit{Fletcher}, this repeal is constitutional because ordinary legislative policies that confer no vested contractual or property rights may always be modified or repealed prospectively by a subsequent legislature (\href{https://digitalcommons.wne.edu/cgi/viewcontent.cgi?article=1919&context=lawreview}{Western New England Law Review}).

\item \textbf{Opposite-Side (Reasonable Exercise of Police Power Over Private Contracts):} During a severe economic depression, a state enacts a temporary, narrowly tailored mortgage moratorium extending the redemption period for defaulting homeowners while requiring them to pay reasonable rental value to lenders. Unlike Georgia's permanent, uncompensated confiscation of vested title in \textit{Fletcher}, modern Contract Clause doctrine upholds reasonable, temporary exercises of the state's police power that address a broad public emergency without permanently destroying the core economic value of contractual rights (as in \textit{Home Building \& Loan Ass'n v. Blaisdell}, 290 U.S. 398 (1934)) (\href{https://supreme.justia.com/cases/federal/us/313/362/}{Justia}).

\item \textbf{Fence-Sitter (Rescission Against Original Corrupt Grantees with Notice):} Suppose a state legislature passes a lucrative mineral-rights lease to a mining syndicate after syndicate executives directly bribe a majority of lawmakers. Before the syndicate transfers any interest to innocent third parties, a newly elected legislature rescinds the lease and files a direct judicial action in equity to cancel the grant for fraud. This is a classic fence-sitter under \textit{Fletcher}: on one hand, Marshall explicitly left open whether a state could void a corrupt grant directly against the original wrongdoers who had notice of the fraud rather than bona fide third-party purchasers; on the other hand, \textit{Fletcher}'s strong reluctance to allow judicial inquiry into legislative motives and its protection of executed public grants create serious doctrinal friction against invalidating a duly enacted statute based on legislative bribery (\href{https://www.oyez.org/cases/1789-1850/10us87}{Oyez}; \href{https://www.fjc.gov/history/cases/cases-that-shaped-the-federal-courts/fletcher-v-peck}{Federal Judicial Center}).
\end{enumerate}

\section{Critique}

\paragraph{Originalist and Textualist Perspectives.}
Originalist and textualist scholars have long questioned Marshall's expansion of Article~I, \S~10. Historical scholarship on the 1787 Convention and the Northwest Ordinance notes that the Contract Clause was primarily aimed at state debtor-relief laws interfering with \textit{private} executory contracts, whereas the Convention's Committee of Style omitted the word ``private'' without recorded debate (\href{https://lawreview.gmu.edu/print__issues/passing-a-complete-code-the-framers-failed-effort-to-nationalize-contract-law/}{George Mason Law Review}). Echoing Justice Johnson's concurrence, textualist critics argue that an executed land grant is a completed conveyance of property rather than an ongoing ``obligation of contracts,'' and that Marshall's invocation of unwritten ``general principles which are common to our free institutions'' improperly blurred positive constitutional text with extratextual natural law (\href{https://scholarship.shu.edu/cgi/viewcontent.cgi?article=2470&context=shlr}{Seton Hall Law Review}; \href{https://scholarship.law.bu.edu/cgi/viewcontent.cgi?article=2045&context=faculty_scholarship}{Boston University Law Scholarship}). Originalists also fault the Court for entertaining a collusive, ``feigned'' lawsuit that strained Article~III's Case-or-Controversy requirement (\href{https://dc.law.utah.edu/cgi/viewcontent.cgi?article=1931&context=ulr}{Utah Law Review}).

\paragraph{Progressive and Doctrinal Perspectives.}
Progressive and realist legal historians criticize \textit{Fletcher v. Peck} as an early milestone of judicial overreach that subordinated democratic self-governance and anti-corruption accountability to speculative commercial and property interests (\href{https://www.fjc.gov/history/cases/cases-that-shaped-the-federal-courts/fletcher-v-peck}{Federal Judicial Center}; \href{https://digitalcommons.wne.edu/cgi/viewcontent.cgi?article=1919&context=lawreview}{Western New England Law Review}). By insulating a notoriously bribed legislature's giveaway of 35 million acres from democratic reversal, the Court signaled that vested corporate and speculative rights could trump popular sovereignty---paving the way for \textit{Dartmouth College} and later nineteenth-century economic substantive due process (\href{https://www.fjc.gov/history/cases/cases-that-shaped-the-federal-courts/fletcher-v-peck}{Federal Judicial Center}). Additionally, federal Indian law scholars critique Marshall's cursory dismissal of indigenous land ownership on the fourth count---treating Georgia as seised in fee simple over lands inhabited by Native nations---as an unexamined doctrinal stepping-stone toward the Doctrine of Discovery later codified in \textit{Johnson v. M'Intosh}, 21 U.S. 543 (1823) (\href{https://www.uclalawreview.org/wp-content/uploads/2019/09/12_52UCLALRev2892004-2005.pdf}{UCLA Law Review}).

\paragraph{Competing Underlying Values.}
\textit{Fletcher v. Peck} exposes an enduring constitutional tension between \textbf{security of vested property rights and commercial stability} (protecting innocent third parties and national markets from retroactive legislative confiscation) and \textbf{democratic accountability and state sovereignty} (empowering voters and successor legislatures to undo corrupt public giveaways). It also pits \textbf{federal judicial supremacy} over state legislation against \textbf{judicial restraint} regarding legislative motive and Article~III justiciability.

\section{Key Quotations}

\begin{quote}
``A contract executed, as well as one which is executory, contains obligations binding on the parties. A grant, in its own nature, amounts to an extinguishment of the right of the grantor, and implies a contract not to reassert that right. A party is, therefore, always estopped by his own grant.''
\end{quote}
\noindent --- Chief Justice John Marshall, Opinion of the Court, 10 U.S. at 137 (\href{https://www.courtlistener.com/opinion/84935/fletcher-v-peck/}{CourtListener}; \href{https://ballotpedia.org/Fletcher_v._Peck}{Ballotpedia}).

\begin{quote}
``Is the clause to be considered as inhibiting the state from impairing the obligation of contracts between two individuals, but as excluding from that inhibition contracts made with itself? The words themselves contain no such distinction. They are general, and are applicable to contracts of every description.''
\end{quote}
\noindent --- Chief Justice John Marshall, Opinion of the Court, 10 U.S. at 137 (\href{https://www.courtlistener.com/opinion/84935/fletcher-v-peck/}{CourtListener}; \href{https://constitutionallawreporter.com/2015/06/16/fletcher-v-peck-contract-clause/}{Constitutional Law Reporter}).

\begin{quote}
``It is, then, the unanimous opinion of the court, that, in this case, the estate having passed into the hands of a purchaser for a valuable consideration, without notice, the state of Georgia was restrained, either by general principles which are common to our free institutions, or by the particular provisions of the constitution of the United States, from passing a law whereby the estate of the plaintiff in the premises so purchased could be constitutionally and legally impaired and rendered null and void.''
\end{quote}
\noindent --- Chief Justice John Marshall, Opinion of the Court, 10 U.S. at 139 (\href{https://www.courtlistener.com/opinion/84935/fletcher-v-peck/}{CourtListener}; \href{https://constitutionallawreporter.com/2015/06/16/fletcher-v-peck-contract-clause/}{Constitutional Law Reporter}).

\begin{quote}
``I do not hesitate to declare that a state does not possess the power of revoking its own grants. But I do it on a general principle, on the reason and nature of things: a principle which will impose laws even on the deity.''
\end{quote}
\noindent --- Justice William Johnson, Jr., Concurring in Part and Dissenting in Part, 10 U.S. at 143 (\href{https://www.courtlistener.com/opinion/84935/fletcher-v-peck/}{CourtListener}; \href{https://caselaw.findlaw.com/court/us-supreme-court/10/87.html}{FindLaw}).

\begin{quote}
``I have been very unwilling to proceed to the decision of this cause at all. It appears to me to bear strong evidence, upon the face of it, of being a mere feigned case. It is our duty to decide on the rights, but not on the speculations of parties.''
\end{quote}
\noindent --- Justice William Johnson, Jr., Concurring in Part and Dissenting in Part, 10 U.S. at 147 (\href{https://www.courtlistener.com/opinion/84935/fletcher-v-peck/}{CourtListener}; \href{https://www.uclalawreview.org/wp-content/uploads/2019/09/12_52UCLALRev2892004-2005.pdf}{UCLA Law Review}).


\end{document}
